\exercisesfor{2}

\begin{enumerate}
\def\labelenumi{\arabic{enumi}.}
\tightlist
\item
  设 g 是具有基本族 \((x_{1},\ldots,x_{n},\ldots)\) 的自由李代数。记 \(g_{n}\) 为由 \(x_{1},\ldots,x_{n}\) 生成的 g 的子代数，并记 \(b_{n}\) 为 \(g_{n}\) 中包含 \(x_{n}\) 的最小理想。
\end{enumerate}

\begin{enumerate}
\def\labelenumi{(\alph{enumi})}
\tightlist
\item
  证明 \(\mathfrak{b}_n\) 是由下列元素生成的 \(\mathfrak{g}_n\) 的子模：
\end{enumerate}

\[
\left(\operatorname{ad} \left(x _ {i _ {1}}\right) \circ \dots \circ \operatorname{ad} \left(x _ {i _ {k}}\right)\right) \left(x _ {n}\right),
\]

其中 \(k \geqslant 0\)，且 \(i_h \leqslant n\) 对所有 \(h\) 都成立。

\begin{enumerate}
\def\labelenumi{(\alph{enumi})}
\setcounter{enumi}{1}
\item
  证明 \(\mathfrak{g}_n = \mathfrak{g}_{n-1} \oplus \mathfrak{b}_n\)；从而推出 \(\mathfrak{g}\) 是各个 \(\mathfrak{b}_n\)（ \(n \geqslant 1\)）的直和。
\item
  利用 (a) 和 (b)，证明 K-模 g 由元素 \([x_{i_{1}}, [x_{i_{2}}, \ldots, [x_{i_{k-1}}, x_{i_{k}}] \ldots]]\) 生成，其中 \(k \geqslant 0\)，且 \(i_{h} \leqslant i_{k}\) 对 h \textless{} k 成立。
\end{enumerate}

¶ 2. 设 X 是至少含有两个元素的可数集，H 是 M(X) 的 Hall 集子集所组成的集合（参见~定义 2）。证明 Card(H) = \(2^{N_0}\)。

\begin{enumerate}
\def\labelenumi{\arabic{enumi}.}
\setcounter{enumi}{2}
\item
  设 \(\mathbf{X}\) 是一个带全序的集合。证明存在一个相对于 \(\mathbf{X}\) 的 Hall 集，使得 \(\mathrm{H} \cap \mathrm{M}^3(\mathbf{X})\) 由乘积 \(z(yx)\)（其中 \(y < x\)、\(y \leqslant z\)）组成，并且 \(\mathrm{H} \cap \mathrm{M}^4(\mathbf{X})\) 由乘积 \(w(z(yx))\)（其中 \(w \geqslant z \geqslant y\)、\(y < x\)）以及乘积 \((ab)(cd)\)（其中 \(a < b\)、\(c < d\)，并且或者 \(a < c\)，或者 \(a = c\) 且 \(b < d\)）组成。
\item
  \begin{enumerate}
  \def\labelenumii{(\alph{enumii})}
  \tightlist
  \item
    证明由表示
  \end{enumerate}
\end{enumerate}

\[
\{x, y; [ x, [ x, y ] ] = [ y, [ x, y ] ] = 0 \}
\]

定义的李代数具有基 \((x,y,[x,y])\)。定义此代数到一个矩阵代数中的嵌入。

\begin{enumerate}
\def\labelenumi{(\alph{enumi})}
\setcounter{enumi}{1}
\tightlist
\item
  对由表示
\end{enumerate}

\[
\{x, y; [ x, [ x, y ] ] - 2 x = [ y, [ x, y ] ] + 2 y = 0 \}.
\]

\begin{enumerate}
\def\labelenumi{(\alph{enumi})}
\setcounter{enumi}{2}
\tightlist
\item
  证明具有表示
\end{enumerate}

\[
\{x, y, z; [ x, y ] - x = [ y, z ] - y = [ z, x ] - z = 0 \}
\]

的李代数化为 0。

\begin{enumerate}
\def\labelenumi{\arabic{enumi}.}
\setcounter{enumi}{4}
\item
  设 \(\mathfrak{g}\) 是自由李代数，\(\mathfrak{r}\) 是 \(\mathfrak{g}\) 的理想。假设 \(\mathfrak{r} = [\mathfrak{g},\mathfrak{r}]\)。证明 \(\mathfrak{r} = \{0\}\)（利用 \(\bigcap_{n\geqslant 0}\mathcal{C}^n\mathfrak{g} = \{0\}\) 这一事实）。
\item
  设 \(\mathbf{E}\) 是一个模。令 \(\mathbf{ME}\) 为 \(\mathbf{E}\) 的自由代数（代数，第三章，§2，习题 13）。回忆 \(\mathrm{ME} = \bigoplus_{n\geqslant 1}\mathrm{E}_n\)，其中
\end{enumerate}

\[
\mathrm{E} _ {1} = \mathrm{E}, \qquad \mathrm{E} _ {n} = \bigoplus_ {p + q = n} \mathrm{E} _ {p} \otimes \mathrm{E} _ {q} \quad \text {if} \quad n \geqslant 2,
\]

且 ME 上的代数结构由典范映射 \(\mathbf{E}_n\otimes \mathbf{E}_m\to \mathbf{E}_{n + m}\) 定义。

\begin{enumerate}
\def\labelenumi{(\alph{enumi})}
\item
  设 J 是 ME 中包含形如 xx 以及 \((xy)z + (yz)x + (zx)y\)（其中 x、y、z 属于 E）的元素的最小双边理想。令 LE = ME/J。证明 J 是 ME 的分次理想；从而得到 LE 的一个分次 \((\mathrm{L}^{n}\mathrm{E})_{n\geqslant1}\)。
\item
  证明 LE 是一个李代数；它称为模 E 的自由李代数。证明对每个李代数 g 和每个线性映射 f: \(E \to g\)，存在唯一一个延拓 f 的李代数同态 F: \(LE \to g\)。
\item
  定义同构 \(\mathbf{E} \to \mathbf{L}^1\mathbf{E}\) 和 \(\wedge^2\mathbf{E} \to \mathbf{L}^2\mathbf{E}\)。证明 \(\mathbf{L}^3\mathbf{E}\) 可与 \(\mathbf{E} \otimes \wedge^2\mathbf{E}\) 对由下列元素生成的子模取商所得的模相识别：
\end{enumerate}

\[
x \otimes (y \wedge z) + y \otimes (z \wedge x) + z \otimes (x \wedge y) \quad \text { for } x, y, z \text { in } E.
\]

\begin{enumerate}
\def\labelenumi{(\alph{enumi})}
\setcounter{enumi}{3}
\item
  证明当 E 是以 X 为基的自由模时，LE 可与第 2 号中定义的自由李代数 \(\mathrm{L}(\mathbf{X})\) 相识别。
\item
  证明 E 到 LE 的包络代数 U(LE) 中的典范映射可以延拓为张量代数 TE 到 U(LE) 的一个同构。
\item
  设 \(\sigma\) 是从 \(\wedge^2\mathrm{E}\) 到 \(\mathsf{T}^2\mathsf{E}\) 的线性映射，满足
\end{enumerate}

\[
\sigma (x \wedge y) = x \otimes y - y \otimes x.
\]

构造一个使 \(\sigma\) 不是单射的模 E。由此推出，对于这个模，LE 到 U(LE) 的典范映射不是单射（参见第一章 §2 的习题 9）。

\begin{enumerate}
\def\labelenumi{\arabic{enumi}.}
\setcounter{enumi}{6}
\tightlist
\item
  设 K 是一个域。令 l 是具有基本族 \((x_{i})_{i\in I}\) 的自由李代数，并令 M 是一个 l-模。
\end{enumerate}

\begin{enumerate}
\def\labelenumi{(\alph{enumi})}
\item
  证明 \(\mathbf{H}^2 (\mathbf{l},\mathbf{M}) = \{0\}\)（参见~第一章 §3 的习题 12。使用命题 1 的推论 2 以及所述习题的 (i) 部分）。
\item
  证明对于 M 中元素的每个族 \((m_{i})_{i\in I}\)，存在唯一的一个 1 次协循环 \(\phi:l\to M\)，满足 \(\phi (x_i) = m_i\)。推出正合列：
\end{enumerate}

\[
0 \rightarrow \mathrm{H} ^ {0} (\mathfrak {l}, \mathrm{M}) \rightarrow \mathrm{M} \rightarrow \mathrm{M} ^ {\mathrm{I}} \rightarrow \mathrm{H} ^ {1} (\mathfrak {l}, \mathrm{M}) \rightarrow 0.
\]

若 I 的基数为有限数 \(n\)，且 M 在 K 上具有有限秩，则

\[
\operatorname{rg} \mathrm{H} ^ {1} (\mathfrak {l}, \mathrm{M}) - \operatorname{rg} \mathrm{H} ^ {0} (\mathfrak {l}, \mathrm{M}) = (n - 1) \operatorname{rg} \mathrm{M}.
\]

\begin{enumerate}
\def\labelenumi{\arabic{enumi}.}
\setcounter{enumi}{7}
\tightlist
\item
  设 K 是一个域。令 l 是具有基本族 \((\overline{x}_{i})_{i\in I}\) 的自由李代数，令 r 是包含于 {[}l, l{]} 中的 l 的理想，并令 g = l/r；用 \(x_{i}\) 表示 \(\overline{x}_{i}\) 在 g 中的像。
\end{enumerate}

证明以下性质等价：

\begin{enumerate}
\def\labelenumi{(\roman{enumi})}
\item
  g 是自由李代数。
\item
  g 以 \((x_{i})\) 为基本族（即~r = \{0\}）。
\item
  对每个 g-模 M，都有 H²(g, M) = \{0\}。
\item
  若 g 在 K 上平凡作用，则 \(H^{2}(g, K) = \{0\}\)。
\end{enumerate}

（蕴含关系 (ii) \(\Rightarrow\) (i) 和 (iii) \(\Rightarrow\) (iv) 显然，而 (i) \(\Rightarrow\) (iii) 由习题 7 得出。为证明 (iv) \(\Rightarrow\) (ii)，只需证明 \(\mathfrak{r} = [\mathfrak{l},\mathfrak{r}]\)，参见~习题 5；否则取 \(\mathfrak{h}\) 为 \(\mathfrak{r}\) 的一个包含 \([\mathfrak{l},\mathfrak{r}]\) 的超平面，并注意到扩张 \(\mathfrak{r} / \mathfrak{h}\to \mathfrak{l} / \mathfrak{h}\to \mathfrak{l} / \mathfrak{r} = \mathfrak{g}\) 是本质的。） 

¶ 9. 设 \(\mathfrak{g} = \bigoplus_{n \geqslant 1} \mathfrak{g}_n\) 是一个分次李代数。若 \(\mathfrak{h}\) 是 \(\mathfrak{g}\) 的分次子代数，且 \(\mathfrak{g} = \mathfrak{h} + [\mathfrak{g}, \mathfrak{g}]\)，证明 \(\mathfrak{h} = \mathfrak{g}\)。若 \((x_i)\) 是 \(\mathfrak{g}\) 的齐次元素族，则推出 \((x_i)\) 是生成族，当且仅当 \(x_i\) 在 \(\mathfrak{g}/[\mathfrak{g}, \mathfrak{g}]\) 中的像生成 \(\mathfrak{g}/[\mathfrak{g}, \mathfrak{g}]\)。

若 K 是 g 在其上平凡作用的域，且 \(\mathrm{H}^{2}(\mathfrak{g},\mathbf{K})=\{0\}\)，证明族 \((x_{i})\) 是基本族，当且仅当 \(x_{i}\) 在 g/{[}g, g{]} 中的像构成这个向量空间的一组基。（使用习题 8。）†

\begin{enumerate}
\def\labelenumi{\arabic{enumi}.}
\setcounter{enumi}{9}
\tightlist
\item
  设 \(t\) 是具有一个含 \(n\) 个元素的有限基本族的自由李代数，并令 \(u\) 是 \(l\) 的满射自同态。证明 \(u\) 是同构。（写出 \(\mathrm{gr}^n(l) = \mathcal{C}^n / \mathcal{C}^{n+1}l\)，并以 \(\mathrm{gr}^n(u)\) 表示 \(\mathrm{gr}^n(l)\) 中由 \(u\) 导出的自同态；注意 \(\mathrm{gr}^n(u)\) 是满射；由于 \(\mathrm{gr}^n(l)\) 是有限秩自由 K-模，推出 \(\mathrm{gr}^n(u)\) 是双射，因为 \(u\) 的核包含于 \(\bigcap_{n} \mathcal{C}^n l\)，而该交为 \(\{0\}\)。）
\end{enumerate}

若 \((y_{1},\ldots ,y_{m})\) 是 l 的生成族，证明 \(m\geqslant n\)，并证明等号成立当且仅当 \((y_{1},\ldots ,y_{m})\) 是基本族。

¶ 11. 设 Y 和 Z 是两个不相交的集合（Y 带有全序），令 \(X = Y \cup Z\) 且 \(l = L(X)\)。令 \(a_{Y}\) 为由 Z 和 {[}l, l{]} 生成的 l 的子模；它是李代数 l 的一个理想。我们拟显式求出 \(a_{Y}\) 的一个基本族。

令 \(\mathbf{M}\) 为 \(\mathbf{Y}\) 上有限支撑的正整数值函数的集合。若 \(m\in \mathbf{M}\)，用 \(\theta^m\) 表示由下式给出的 l 的自同态：

\[
\theta^ {m} = \prod_ {y \in Y} (a d y) ^ {m (y)},
\]

其中乘积按 Y 上给定的序关系进行。令 \(\mathfrak{S}\) 表示 \(\mathbf{Y} \times \mathbf{M}\) 的子集，它由满足下述条件的有序对 \((u, m)\) 构成：存在 \(y \in \mathbf{Y}\) 使得 \(y > u\) 且 \(m(y) \geqslant 1\)。证明元素

\[
\theta^ {m} (u) \quad \text { for } \quad (u, m) \in \mathfrak {S} \quad \text { and } \quad \theta^ {m} (z) \quad \text { for } \quad (z, m) \in Z \times M
\]

构成李代数 \(\mathfrak{a}_{\mathbf{Y}}\) 的一个基本族。

（将问题归结为 Y 有限的情形，并按 Card(Y) 归纳。将命题 10 的推论应用于 Y 的最大元 y，并将归纳假设应用于 \(Y = \{y\}\)。）

\begin{enumerate}
\def\labelenumi{\arabic{enumi}.}
\setcounter{enumi}{11}
\tightlist
\item
  设 \(\mathbf{X} = \{x, y\}\) 是一个含两个元素的集合。证明 \(\mathbf{L}(\mathbf{X})\) 的导出代数是自由李代数，并以元素 \(((\operatorname{ad}y)^q \circ (\operatorname{ad}x)^p)(y)\) 为基本族，其中 \(p \geqslant 1\)、\(q \geqslant 0\)。（使用习题 11。）
\end{enumerate}

¶ 13.（本习题中假设 K 上的每个投射模都是自由的；例如当 K 为主理想整环时即如此。）

令 \(\mathfrak{l} = \sum_{n=1}^{\infty} \mathfrak{l}_n\) 是一个具有由齐次元素组成的基本族 B 的分次李代数，并令 \(\mathfrak{h}\) 是 \(\mathfrak{l}\) 的一个分次李子代数；作为模看，它是 \(\mathfrak{l}\) 的一个直因子。

\begin{enumerate}
\def\labelenumi{(\alph{enumi})}
\tightlist
\item
  对所有 \(i \geqslant 0\)，令 \(\mathfrak{l}^{(i)}\) 为 \(\mathfrak{l}\) 的分次子代数，满足
\end{enumerate}

\[
\begin{array}{l l} \mathfrak {l} _ {j} ^ {(i)} = \mathfrak {h} _ {j} & \text { if } j \leqslant i. \\ \mathfrak {l} _ {j} ^ {(i)} = \mathfrak {l} _ {j} & \text { if } j > i. \end{array}
\]

于是 \(l = \ell^{(0)} \supset \ell^{(1)} \supset \cdots \supset \mathfrak{h}\)。证明通过对 \(i\) 归纳，存在一个记作 \(B^{(i)}\) 的 \(\ell^{(i)}\) 的基本族，且由齐次元素组成，使得 \(B^{(i-1)}\) 中以及 \(B^{(i)}\) 中次数 \(< i - 1\) 的元素相同。（假设 \(B^{(i-1)}\) 已构造。令 \(m_i\) 为 \(l_i = l_i^{(i-2)}\) 与由 \(b_j\)（其中 \(j < i\)）生成的子代数的交，并令 \(b_i\) 为 \(l_i\) 中由 \(B^{(i-1)}\) 中次数为 \(i\) 的元素生成的子模。归纳假设蕴含 \(l_i = m_i \oplus b_i\)。由于 \(m_i \subset b_i\)，可以将 \(b_i\) 分解为 \(b_i = y_i \oplus b_i\)，从而 \(b_i = m_i \oplus b_i\)；由关于 K 的假设，模 \(y_i\) 和 \(b_i\) 是自由的；

必要时改变 \(\mathrm{B}^{(i-1)}\) 中次数为 i 的元素，可以假设 \(b_{i}\) 由 \(\mathrm{B}^{(i-1)}\) 的一个子集生成。将习题 11 应用于 \(\mathfrak{l}^{(i-1)}\) 及其理想 \(\mathfrak{l}^{(i)}\)，得到一个记作 \(\mathrm{B}^{(i)}\) 的 \(\mathfrak{l}^{(i)}\) 的基本族，具有所需性质。）

\begin{enumerate}
\def\labelenumi{(\alph{enumi})}
\setcounter{enumi}{1}
\tightlist
\item
  由 (a) 推出 \(\mathfrak{h}\) 具有由齐次元素组成的基本族。
\end{enumerate}

¶ 14. 设 K 是一个域。令 X 是一个集合，x、y 是 L(X) 中关于 K 线性无关的两个元素。证明族 \((x,y)\) 在李代数 L(X) 中是自由的。（令 \(x_{p}\)（相应地 \(y_{q}\)）为 x（相应地 y）的最高次数非零齐次分量。必要时将 y 的一个倍数加到 x 上，可以假设 \(x_{p}\) 和 \(y_{q}\) 线性无关。由 \(x_{p}\) 和 \(y_{q}\) 生成的 L(X) 的子代数是分次的，因而是自由的，参见~习题 13。推出 \((x_{p},y_{q})\) 是自由族，再由此过渡到 \((x,y)\)。）

¶ 15. 设 \(\mathrm{X} = \{x, y\}\) 是一个含两个元素的集合，且 \(\sigma\) 是李代数 \(\mathrm{L}(\mathrm{X})\) 的一个自同构。证明若 \(\mathrm{K}\) 是域，则 \(\sigma\) 保持 \(\mathrm{L}(\mathrm{X})\) 的分次（将习题 14 应用于 \(\sigma(x)\) 和 \(\sigma(y)\) 的最高次数非零齐次分量）；从而推出 \(\operatorname{Aut}(\mathrm{L}(\mathrm{X}))\) 同构于 \(\mathbf{GL}(2, \mathrm{K})\)。将这些结果推广到环 \(\mathrm{K}\) 没有幂零元 \(\neq 0\) 的情形。当 \(\mathrm{K}\) 含有平方为零的元素 \(\varepsilon \neq 0\) 时，证明 \(x \mapsto x, y \mapsto y + \varepsilon[x, y]\) 可以延拓为 \(\mathrm{L}(\mathrm{X})\) 的一个不保持 \(\mathrm{L}(\mathrm{X})\) 的分次的自同构。

¶ 16. 设 K 是一个 Q-代数。若 g 是一个李代数，用 U(g) 表示其包络代数，用 σ 表示 g 到 U(g) 的典范映射，并用 S(g) 表示 K-模 g 的对称代数。存在唯一的线性映射

\[
\eta (\mathfrak {g}): \mathrm{S} (\mathfrak {g}) \rightarrow \mathrm{U} (\mathfrak {g})
\]

满足 \(\eta (\mathfrak{g})(x^n) = \sigma (x)^n\) 对所有 \(x\in \mathfrak{g}\) 和所有 \(n\in \mathbf{N}\) 均成立。于是

\[
\eta (\mathfrak {g}) (x _ {1} \dots x _ {n}) = \frac {1}{n !} \sum_ {s \in \mathfrak {S} _ {n}} \sigma (x _ {s (1)} \dots x _ {s (n)}), \quad x _ {i} \in \mathfrak {g}.
\]

我们拟证明 \(\eta(g)\) 是双射。

\begin{enumerate}
\def\labelenumi{(\alph{enumi})}
\item
  证明 \(\eta(g)\) 是满射；当 K-模 \(g\) 是自由模时，证明它是双射（使用庞加莱–伯克霍夫–维特定理）。
\item
  令 \(\mathfrak{h}\) 是 \(\mathfrak{g}\) 的一个理想。令 \(\mathfrak{s}_{\mathfrak{h}}\)（相应地，\(\mathfrak{u}_{\mathfrak{h}}\)）表示 \(S(\mathfrak{g})\)（相应地，\(U(\mathfrak{g})\) 的双边理想）中由 \(\mathfrak{h}\)（相应地，\(\sigma(\mathfrak{h})\)）生成的理想。下图
\end{enumerate}

\[
\begin{array}{c} 0 \to s _ {\mathfrak {h}} \to S (g) \to S (g / \mathfrak {h}) \to 0 \\ \downarrow^ {\eta (g)} \quad \downarrow^ {\eta (g / \mathfrak {h})} \\ 0 \to u _ {\mathfrak {h}} \to U (g) \to U (g / \mathfrak {h}) \to 0 \end{array}
\]

图交换，且其各行正合。由此推出 \(\tilde{u}_{\mathfrak{h}}\)，即 \(s_{\mathfrak{h}}\) 在 \(\eta(g)\) 下的像，包含于 \(u_{\mathfrak{h}}\) 中；并且 \(\tilde{u}_{\mathfrak{h}} = u_{\mathfrak{h}}\)，只要 \(\eta(g)\) 和 \(\eta(g/h)\) 是单射（特别地，当模 \(g\) 和 \(g/h\) 是自由模时）。

\begin{enumerate}
\def\labelenumi{(\alph{enumi})}
\setcounter{enumi}{2}
\item
  取 \(\mathfrak{g}\) 为具有基本族 \((\mathbf{X}_1, \ldots, \mathbf{X}_n, \mathrm{H})\) 的自由李代数，其中 \(n \geqslant 0\)，并令 \(\mathfrak{h}\) 为由 H 生成的 \(\mathfrak{g}\) 的理想。证明 \(\mathfrak{g}\) 和 \(\mathfrak{g} / \mathfrak{h}\) 是自由模。由此推出在该情形下 \(\tilde{\mathfrak{u}}_{\mathfrak{h}} = \mathfrak{u}_{\mathfrak{h}}\)。特别地，存在 \(x \in \mathfrak{s}_{\mathfrak{h}}\)，使得 \((\eta(\mathfrak{g}))(x) = \sigma(\mathbf{X}_1) \ldots \sigma(\mathbf{X}_n) \sigma(\mathrm{H})\)。
\item
  回到一般情形。证明 \(\mathfrak{u}_{\mathfrak{h}}\) 作为 \(\mathrm{K}\)-模由形如 \(\sigma(x_1) \ldots \sigma(x_n)\sigma(h)\) 的元素生成，其中 \(n \geqslant 0\)、\(x_i \in \mathfrak{g}\) 且 \(h \in \mathfrak{h}\)（注意 \(\mathfrak{u}_{\mathfrak{h}}\) 与由 \(\mathfrak{h}\) 生成的左理想重合）。利用 (c) 证明这种元素属于 \(\tilde{\mathfrak{u}}_{\mathfrak{h}}\)（使用从自由李代数到 \(\mathfrak{g}\) 的适当同态）。推出 \(\tilde{\mathfrak{u}}_{\mathfrak{h}} = \mathfrak{u}_{\mathfrak{h}}\)。
\item
  证明若 \(\eta(g)\) 是双射，则 \(\eta(g/h)\) 也是双射。最后推出，对每个李代数 \(g, \eta(g)\) 都是双射（将 \(g\) 写成一个自由李代数的商），并且典范同态
\end{enumerate}

\[
\omega : \mathsf {S} (\mathfrak {g}) \to \operatorname{gr} \mathrm{U} (\mathfrak {g}) \quad (\text { cf.   Chapter   I,   } \S 2, \text {   no.   } 6)
\]

是同构（Q-代数上李代数的“庞加莱–伯克霍夫–维特定理”）。

